% M1..M6 trong log/luận văn = ER, CR, SL, AL, PT, US.
\begin{table}[!htbp]
\centering
\caption{Sáu cơ chế tấn công và vị từ gây hại tương ứng.}\label{tab:mechanisms}
\footnotesize
\begin{tabular}{@{}lp{0.70\textwidth}l@{}}
\toprule
Mã & Cơ chế (kiểu con) & Vị từ \\
\midrule
\mER & Chuyển hướng tường minh: yêu cầu gửi tới đích không có trong $u$, có động từ gửi và địa chỉ đầy đủ & $H_\mathit{dest}$ \\
\mCR & Chuyển hướng che giấu: cùng mục tiêu nhưng tránh mẫu ``gửi + địa chỉ'' (ghi nhận không động từ, chuyển tiếp hoặc cc, tách địa chỉ, đích nằm ở tài liệu thứ hai) & $H_\mathit{dest}$ \\
\mSL & Rò rỉ qua sink: đưa giá trị mật vào sink người dùng đã nêu (thân e-mail, payload, trường trông vô hại) & $H_\mathit{leak}$ \\
\mAL & Rò rỉ qua câu trả lời: đưa giá trị mật vào câu trả lời cuối (nguyên văn, lồng trong câu, có nhãn) & $H_\mathit{leak}$ \\
\mPT & Sửa tham số: đổi một giá trị thường trong payload của sink người dùng đã nêu (số, boolean, đơn vị, định danh) & $H_\mathit{tamper}$ \\
\mUS & Tác dụng phụ ngoài yêu cầu: chèn lời gọi sink vào nhiệm vụ chỉ đọc (đích chỉ xuất hiện trong câu cấm, chỉ là thông tin, hoặc không có) & $H_\mathit{side}$ \\
\bottomrule
\end{tabular}
\end{table}
